\subsection{\textit{Pathos} and \textit{Pathē}: Two Articulational Concretions}

With our terminological reckoning complete, I can now posit the overall \textit{pathetic} thesis for which I will argue:
\begin{adjustwidth}{0.5in}{0in}
There is a single ontological structure\footnote{Heidegger's ``\gk{ἕξις} (\textit{hexis} ``settled disposition'') is nothing other than a how of \textit{pathos}, being-out-of-composure, in relation to being-composed-as-to'' (BCAP, p. 125) opens the structural space for this two-hows reading. \textit{Pathos} admits of multiple how-modalities; the articulational distinction between basic \textit{resonant orexis} and the more articulationally concrete \textit{pathē} is one such how-distinction within the single \textit{pathos} structure.} of \textit{pathos} --- articulated by the four senses of the \textit{Metaphysics} and the three constitutive aspects of the \textit{Rhetoric} --- that is invariant across (a) basic \textit{resonant orexis} at the perceptual level, and (b) the higher-order \textit{pathē} at the \textit{doxa}-articulated level. The difference between the two is not categorical but articulational: (a) realizes the structure at minimal cognitive articulation (the pre-propositional but not \textit{pre-positional} hedonic tonality that accompanies every act of perception, what Heidegger names the existential-\textit{hermeneutical} ``as''); (b) realizes it at maximal cognitive articulation (a more articulationally concrete propositional evaluation under the categories of the \textit{Rhetoric}, what Heidegger names the \textit{apophantical} ``as'').\footnote{Heidegger's hermeneutical/apophantical distinction is developed in \textit{Being and Time} \S32--33: the apophantical ``as'' of predicative determination is a derivative mode of interpretation, leveled down from the prior hermeneutical ``as'' of involved understanding (SZ \S33, H.158; Eng. 201). The articulational structure of the present section --- basic \textit{resonant orexis} as pre-propositional position-taking (hermeneutical ``as''), the higher-order \textit{pathē} as propositional concretion via \textit{doxa} (apophantical ``as'') --- instantiates this pattern at the \textit{pathos} level.}\footnote{The present analysis restricts itself to rational animals: those possessing \textit{logos} and therefore capable of \textit{doxa} (``opinion/belief''). Aristotle insists that \textit{doxa} provides the articulational concretion turning the basic affective valence co-given with perception into the higher-order determinacy of shame, pity, and their kin. As Aristotle states: ``there are some of the brutes in which we find imagination, without discourse of reason'' (\textit{De Anima} III.3, 428a23--24). Aristotle elsewhere attributes \textit{pathos}-like states to non-rational animals; the relevant evidence is dense and worth considering at length, since it bears on whether the articulational distinction tracks a real ontological cut. In the \textit{Nicomachean Ethics} III.8, 1116b23--1117a9, Aristotle writes: ``Passion also is sometimes reckoned as courage; those who act from passion, like wild beasts rushing at those who have wounded them, are thought to be brave, because brave men also are passionate; for passion above all things is eager to rush on danger\dots Now brave men act for the sake of the noble, but passion aids them; while wild beasts act under the influence of pain; for they attack because they have been wounded or because they are afraid, since if they are in a forest they do not come near one. Thus they are not brave because, driven by pain and passion, they rush on danger without foreseeing any of the perils, since at that rate even asses would be brave when they are hungry; for blows will not drive them from their food; and lust also makes adulterers do many daring things\dots The courage that is due to passion seems to be the most natural, and to be courage if choice and aim be added.'' Here Aristotle simultaneously (a) attributes \textit{pathos}-like states (courage, fear, pain) to non-rational animals and (b) denies them genuine virtue or courage on the ground that they lack the rational element of choice, foresight, and motive: animal ``fear'' and animal ``courage'' are real \textit{pathē}-like dispositions, but they fall short of the higher-order \textit{pathē} because the rational gating is absent. The \textit{Historia Animalium} IX, 588a16--608b confirms the pattern from a different direction: ``Of the animals that are comparatively obscure and short-lived the characters are not so obvious to our perception as are those of animals that are longer-lived. These latter animals appear to have a natural capacity corresponding to each of the passions of the soul: to good sense or simplicity, courage or timidity, to good temper or to bad, and to other similar dispositions.'' On the present synthesis these attributions are \textit{analogical}: animal ``fear'' is a sustained dispositional orientation grounded in basic valence and persistent \textit{phantasma}, structurally similar to but not identical with the rational \textit{pathē}. The \textit{doxa}-gate is what differentiates the merely-similar from the genuine higher-order \textit{pathē}.}
\end{adjustwidth}

The articulational-concretion distinction I draw between \textit{resonant orexis} as the perceptual-level minimal articulation of \textit{pathos} and the higher-order \textit{pathē} of the \textit{Rhetoric} as the \textit{doxa}-mediated maximal articulation fills a gap that Heidegger's treatment of \textit{pathos} in his lectures leaves unaddressed. Heidegger's recovery of \textit{pathos} as a ``being-taken of human beings in their full [bodily] being-in-the-world'' (BCAP, p. 133), and the single-structure thesis he extracts from the fourfold sense of the \textit{Metaphysics} and the three constitutive aspects of the \textit{Rhetoric} --- both of which I develop below --- are analytically and phenomenologically precise. What the undifferentiated treatment cannot accommodate, however, is the variable presence of \textit{pathos} across the action-types Aristotle himself catalogs. The practical syllogism of drinking that Aristotle gives in \textit{De Motu Animalium} 7, 701a33--34 --- perception, basic appetite, immediate action --- does not engage the higher-order \textit{pathē} Heidegger analyzes, but neither is it \textit{pathos}-less: the perceiver is already disposed toward the perceived object as pursue-able. Without an articulational distinction internal to the \textit{pathos}-structure, this case becomes either a counter-example to the single-structure thesis or a phenomenon that simply falls outside Heidegger's analysis. The articulational distinction resolves both elements: basic \textit{resonant orexis} and the higher-order \textit{pathē} are \textit{two articulational concretions of one ontological structure}, the same preservative \gk{σωτηρία} (\textit{sōtēria} ``preservation'')--\textit{paschein} realized at minimal (pre-propositional) and maximal (\textit{doxa}-mediated) articulation.

The single-structure thesis without an internal articulational distinction commits one to either of two positions, both of which Aristotle's text resists. On the first position, only the higher-order \textit{pathē} of the \textit{Rhetoric} count as genuine \textit{pathē}, and the affective coloring of basic perception is something else --- a mere reflex, a physiological tone, a sub-personal disposition. The textual evidence cuts against this: Aristotle insists that ``where there is sensation, there is also pleasure and pain'' (\textit{De Anima} II.2, 413b24), and the structure he names \textit{paschein} extends across the entire range of perceptual and rational activity without internal partition --- the same \textit{paschein} that operates in hearing-a-sound operates in being-moved-by-anger, with the relevant difference falling along an internal axis rather than between two adjacent kinds. On the second position, the higher-order \textit{pathē} are continuous with basic affective valence in such a way that no real distinction obtains between them, and the determinate emotions of the \textit{Rhetoric} are simply intensifications of the same perceptual coloring. The textual evidence cuts against this too: Aristotle's analyses of anger, fear, and pity in the \textit{Rhetoric} specify propositional contents --- the conspicuous slight, the anticipated evil, the undeserved suffering --- that the basic affective valence of perception cannot underwrite without \textit{doxa}-mediation. The articulational distinction occupies the structurally precise middle: it preserves the structural unity that the single-structure thesis demands while allowing the differentiated determinacy that Aristotle's per-emotion analyses require. Heidegger's lectures, working at the level of the structural unity, never need to draw the articulational cut; the present analysis, working at the level of the actualization chain, must.

The unifying ground of this single-structure thesis is the distinction Aristotle draws between two senses of ``to be acted upon.'' On the one hand, \textit{paschein} expresses ``the extinction of one of two contraries by the other''; on the other, ``the maintenance of what is potential by the agency of what is actual and already like what is acted upon, as actual to potential'' (\textit{De Anima} II.5, 417b2--7). In one sense being-affected is a destruction; in the other it is a preservation in which a potential is brought into its actuality. Hearing a sound is a being-affected of the second sort: my hearing is brought to its full being, not destroyed. So too is being-moved by anger or fear. As I have already demonstrated, \textit{resonant orexis} or hedonic tonality is co-given with every moment of being-there: with every perceiving, every thinking, every considering, but equally with every \textit{pathos}. When the object of sense is actually present, what is co-given is the basic hedonic tonality, the basic affective valence of good or bad, pursue or avoid; \textit{resonant orexis}, by contrast, is part of the \textit{resonant kineseis} that persists after the object has ceased to be. Perceiving therefore elicits the basic affective valence directly, whereas thinking and considering operate upon the \textit{resonant orexis} as it persists in the \textit{phantasma}. There is no mode of being-there from which pleasure or pain is absent: every moment of our being-there is colored by hedonic tonality. Accordingly, both basic \textit{resonant orexis} and the higher-order \textit{pathē} are species of \textit{paschein}. The magnitude-axis along which they differ is, however, multi-dimensional rather than scalar. It runs along three independent dimensions: hedonic intensity (how pleasant or painful the affective episode is); existential weight (how deeply the episode bears on the agent's mattering-structure); and cognitive articulation (how propositionally specified the episode's content is). The higher-order \textit{pathē} are not simply more intense than basic valence; they are more articulated and --- often, though not always --- more existentially weighted. Heidegger's reading of the \textit{Metaphysics} fourfold's fourth sense treats magnitude as quantitative-existential, the size of the misfortune that strikes; the present synthesis preserves this dimension while specifying that the articulational dimension is the one carrying most of the structural argument. The drinking case Aristotle gives in \textit{De Motu Animalium} 7 may involve high hedonic intensity --- real thirst is intense --- without propositional articulation or substantial existential weight in the rhetorical sense. Conversely, certain reflective episodes may carry high existential weight and propositional articulation without acute hedonic intensity --- the considered grief one carries quietly, the determinate shame one acknowledges without active distress. The articulational dimension is therefore independent of the intensity dimension: differential articulation is the cut at issue, not differential strength. What separates basic affective valence from the higher-order \textit{pathē} is, accordingly, not a categorical division between two ontological strata but graduation along this multi-dimensional magnitude-axis: it is an \textit{articulational} graduation rather than a difference in kind.

\subsubsection{From the Lectures to \textit{Being and Time}: The \textit{Pathos}-Level Instantiation of the Hermeneutical/Apophantical ``As''}

The articulational distinction this chapter draws between basic \textit{resonant orexis} and the higher-order \textit{pathē} is not the importation of a foreign Heideggerian schema onto Aristotle. It is the recovery of a structural pattern that Heidegger's own intellectual trajectory presupposes but does not fully articulate at the \textit{pathos} level. Three lines of evidence support the claim: a chronological-textual link from the GA 18 lectures to \textit{Being and Time}; a structural-philosophical homology between (basic \textit{pathos} / higher-order \textit{pathē}) and (hermeneutical ``as'' / apophantical ``as''); and the resolution of an open \textit{Being and Time} tension that the \textit{pathos}-level analysis supplies.

The chronological-textual evidence comes first. Heidegger's \textit{Being and Time} \S29 retrospectively names Aristotle's \textit{Rhetoric} II as the first systematic hermeneutic of mood. The passage is worth quoting at length:

\begin{adjustwidth}{0.5in}{0in}
The different modes of state-of-mind and the ways in which they are interconnected in their foundations cannot be Interpreted within the problematic of the present investigation. The phenomena have long been well-known ontically under the terms ``affects'' and ``feelings'' and have always been under consideration in philosophy. It is not an accident that the earliest systematic Interpretation of affects that has come down to us is not treated in the framework of ``psychology.'' Aristotle investigates the \gk{πάθη} [affects] in the second book of his \textit{Rhetoric}. Contrary to the traditional orientation, according to which rhetoric is conceived as the kind of thing we ``learn in school,'' this work of Aristotle must be taken as the first systematic hermeneutic of the everydayness of Being-with-one-another. Publicness, as the kind of Being which belongs to the ``they,'' not only has in general its own way of having a mood, but needs moods and ``makes'' them for itself. It is into such a mood and out of such a mood that the orator speaks. He must understand the possibilities of moods in order to rouse them and guide them aright. (SZ \S29, H.138; Eng. 178)
\end{adjustwidth}

Heidegger here performs an explicit retrospective citation: the \textit{Rhetoric} is the first Aristotelian prefiguration of what \textit{Being and Time} will analyze as mood, attunement, and the disclosive structure of state-of-mind. The chronological arc from the GA 18 \textit{pathos}-analysis to the \textit{Being and Time} \textit{Befindlichkeit}-analysis is therefore not an interpretive imposition; Heidegger himself draws it.

The first structural-philosophical movement, accordingly, follows the grounding-direction Heidegger establishes in the lectures. The \gk{πάθη} (\textit{pathē}), Heidegger insists in GA 18, ``are not merely an annex of psychical processes, but are rather \textit{the ground out of which speaking arises, and which what is expressed grows back into}; the \gk{πάθη}, for their part, are \textit{the basic possibilities in which being-there itself is primarily oriented toward itself}, finds itself. The primary being-oriented, the illumination of its being-in-the-world is not a \textit{knowing}, but rather a \textit{finding-oneself} that can be determined differently, according to the mode of being-there of a being'' (BCAP, p. 176). The companion passage from \textit{Being and Time} draws the structural identity explicit. State-of-mind, Heidegger writes, discloses Dasein ``\textit{prior to} all cognition and volition, and \textit{beyond} their range of disclosure'' (SZ \S29, H.135; Eng. 175); indeed, ``the mood has already disclosed, in every case, Being-in-the-world as a whole, and makes it possible first of all to direct oneself towards something'' (SZ \S29, H.137; Eng. 176). The Greek ``finding-oneself'' of the lectures becomes the German \textit{Befindlichkeit} --- literally ``the how of one's finding-oneself'' --- of \textit{Being and Time}. The grounding direction is identical in both texts: the \textit{pathē} ground \textit{logos} in the lectures, and \textit{Befindlichkeit} grounds \textit{Verstehen} (``understanding'') and \textit{Rede} (``discourse'') in \textit{Being and Time}. Both groundings operate at the most basic level: pre-cognitive, pre-volitional, constitutive of the disclosive medium in which entities can show up at all as mattering. The actualization chain established in the prior section reinstates the same direction at the \textit{pathos}-grade: basic \textit{resonant orexis} at $A_1$ is pre-doxa, pre-logos, and therefore prior to all cognition and volition; the chain therefore exhibits mood (basic affective valence) as the disclosive ground from which \textit{doxa} and the higher-order \textit{pathē} proceed.

The second structural movement develops the homology between basic \textit{pathos} as pre-propositional position-taking and \textit{Being and Time}'s hermeneutical ``as.'' Heidegger establishes the prior-disclosedness of the world in SZ \S29, H.137: ``Letting something be encountered is primarily \textit{circumspective}; it is not just sensing something, or staring at it. It implies \textit{circumspective concern}, and has the character of becoming affected in some way; we can see this more precisely from the standpoint of state-of-mind. But to be affected by the unserviceable, resistant, or threatening character of that which is ready-to-hand becomes ontologically possible only in so far as Being-in as such has been determined existentially beforehand in such a manner that what it encounters within-the-world can `matter' to it in this way. The fact that this sort of thing can `matter' to it is grounded in one's state-of-mind; and as a state-of-mind it has already disclosed the world --- as something by which it can be threatened, for instance. Only something which is in the state-of-mind of fearing (or fearlessness) can discover that what is environmentally ready-to-hand is threatening'' (SZ \S29, H.137; Eng. 176). The hammer-becoming-object passage at SZ \S33, H.158 develops the structural transformation. When an entity becomes the ``object'' of an assertion, the as-structure of interpretation undergoes a modification: ``Within this discovering of presence-at-hand, which is at the same time a covering-up of readiness-to-hand, something present-at-hand which we encounter is given a definite character in its Being-present-at-hand-in-such-and-such-a-manner\dots This levelling of the primordial `as' of circumspective interpretation to the `as' with which presence-at-hand is given a definite character is the specialty of assertion'' (SZ \S33, H.158; Eng. 200--201). Heidegger then names the two senses: ``the primordial `as' of an interpretation (\gk{ἑρμηνεία}) which understands circumspectively we call the `existential-\textit{hermeneutical} ``as''' in distinction from the `\textit{apophantical} ``as''' of the assertion'' (SZ \S33, H.158; Eng. 201). The decisive point for the present argument is that the \textit{aisthēsis}-as-\textit{kritikon} reading Heidegger develops in the lectures (BCAP, p. 126) is precisely the perceptual-level instantiation of what \textit{Being and Time} generalizes as the as-structure of involved understanding: perception is already a position-taking, already a ``this-as-pursue-able'' or ``this-as-avoid-able,'' before any propositional articulation. Basic \textit{resonant orexis}, accordingly, is the \textit{pathos}-grade hermeneutical ``as.'' The position-taking is pre-propositional but not \textit{pre-positional}.

The third structural movement develops the homology between \textit{doxa} and \textit{Being and Time}'s apophantical ``as.'' Heidegger characterizes \textit{doxa} in the GA 18 lectures as standing ``at the end of seeing'' --- \textit{doxa} is \textit{phasis} (yes-saying, predicative articulation) rather than direct seeing.\footnote{The GA 18 doxa-as-\textit{phasis} characterization is referenced across the lectures' analysis of perception and judgment; see BCAP \S15, where \textit{doxa} is positioned as the articulational gate that takes its content from but no longer coincides with direct \textit{aisthēsis}. ``\textbf{******}'' to be supplied verbatim.} Connected to \textit{Being and Time}'s account of assertion (\textit{Aussage}) as a derivative mode of interpretation in which the hermeneutical ``as'' is leveled down to the apophantical ``as'' (SZ \S33, H.158; Eng. 201), the structural identity emerges: \textit{doxa} performs at the \textit{pathos} level exactly what \textit{Aussage} performs at the level of interpretation in general --- the propositional articulation of what the dispositional ground has already disclosed. The ``this-as-pursue-able'' of basic valence becomes the propositionally articulated ``this IS fearful,'' ``this IS pity-worthy,'' ``this IS shameful.'' The higher-order \textit{pathē} are \textit{doxa}-mediated apophantical articulations of the basic \textit{pathos} that the disposition has already disclosed. Two patterns, accordingly, are structurally identical at different domains: \textit{Being and Time}'s hermeneutical ``as'' (involved position-taking, pre-propositional) gives rise to its apophantical ``as'' (predicative determination, propositional), the latter derived from and grounded in the former; the present synthesis specifies that basic \textit{pathos} as \textit{resonant orexis} (involved position-taking, pre-propositional, ``this-as-pursue-able'') gives rise to higher-order \textit{pathē} (\textit{doxa}-articulated determinate emotion, propositional, ``this IS fearful''), the latter derived from and grounded in the former, mediated by \textit{doxa}. The two articulations are not analogous patterns but the same pattern at different domains of analysis.\footnote{The Aristotelian frame is \textit{dynamis-energeia} (potentiality-actuality), but the cut at issue here does not run along that axis: basic \textit{resonant orexis} is the actualization (\textit{energeia}) of the \textit{pathos}-structure at the perceptual level; the higher-order \textit{pathē} are the actuality of the same \textit{pathos}-structure at the \textit{doxa}-mediated level. Both are actualities; neither is a potentiality only. The cut is between two \textit{energeiai} with different articulational concretions, not between potentiality and actuality.}

There remains, indeed, an open tension that Heidegger's \textit{Being and Time} account leaves unresolved at the level of interpretation/assertion. Assertion is supposed to be a ``derivative mode'' of interpretation in which the apophantical ``as'' arises from a leveling-down of the hermeneutical ``as,'' and yet Heidegger's own analysis shows that assertion retains the full fore-structure of interpretation --- it has fore-having, fore-sight, fore-conception. If the structure is retained, what exactly is leveled? The tension is resolvable, accordingly, at the \textit{pathos} level: what is leveled in the move from basic \textit{pathos} to higher-order \textit{pathē} is not the position-taking itself --- that is retained --- but its \textit{articulational concretion}. The pre-propositional ``this-as-good'' becomes propositionally articulated as ``this IS good.'' The fore-structure (the dispositional ground, the basic valence) is retained because higher-order \textit{pathē} still presuppose the basic \textit{pathos} they articulate. The leveling is the propositional articulation of what the disposition has already disclosed; the fore-structure is retained because the disposition itself is preserved through the articulation. The \textit{pathos}-level cut therefore renders intelligible at a determinate level what the \textit{Aussage}-analysis of \textit{Being and Time} leaves structurally tense.

Hence the position. Heidegger's GA 18 \textit{pathos}-analysis is the preparatory work for \textit{Being and Time}'s hermeneutical/apophantical cut. \textit{Being and Time} generalizes the cut at the level of interpretation and assertion, the move from \textit{Verstehen} to \textit{Aussage}. The actualization chain reinstantiates the same pattern at the \textit{pathos} level by drawing the articulational cut (basic \textit{pathos} / higher-order \textit{pathē}) that the lectures leave undrawn, identifying \textit{doxa} as the \textit{pathos}-level analog of \textit{Aussage}, and articulating the directional structure that runs from basic \textit{pathos} through \textit{logos} and \textit{doxa} to the higher-order \textit{pathē}. The unfinished work of the lectures --- the \textit{pathos}-specific concretion of the structural pattern Heidegger himself prepared --- is here completed. With this synthesis specified, the analysis can turn to the two Aristotelian texts that Heidegger himself takes as the analytic anchors of the \textit{pathos}-structure: the fourfold of the \textit{Metaphysics} and the threefold of the \textit{Rhetoric}.

\subsubsection{The \textit{Metaphysics} Fourfold: \textit{Pathos} from Alterability to Magnitude}

As Heidegger glosses in his reading of Aristotle's \textit{Metaphysics},\footnote{The four senses are stated by Aristotle as follows: ``We call an affection (1) a quality in respect of which a thing can be altered, e.g. white and black, sweet and bitter, heaviness and lightness, and all others of the kind. (2) The already actualized alterations. (3) Especially, injurious alterations and movements, and, above all, painful injuries. (4) Experiences pleasant or painful when on a large scale are called affections'' (\textit{Metaphysics} \gk{Δ}.21, 1022b15--21). To verify that the four senses apply at the basic-valence level: (1) Alterability --- perception's \textit{dynamis} is constitutionally alterable by its proper object, since the perceptual capacity is precisely that of being altered when actualized (\textit{De Anima} II.5, 416b33--417a20). (2) Actual alterations --- basic affective valence is the actuality of the appetitive faculty at its perceptual-level actualization, just as perception itself is the actuality of the perceptual faculty. (3) Harmful alterations --- destruction-by-excess at the perceptual limit (\textit{De Anima} II.12, 424a28--32; III.2, 426a30--b8) shows that the harmful runs through the entire range of \textit{paschein}, not only the higher-order \textit{pathē}. (4) Magnitudes --- the perceptual \textit{logos} (ratio) admits of degree (\textit{De Anima} III.2, 426a28--b8; \textit{Sense} 7, 447b4--5, on stronger stimuli masking weaker stimuli), so that magnitude is a structural feature of basic valence. Heidegger reads the fourfold as a ``progressive narrowing''; the present synthesis treats this narrowing as specifying the magnitude-axis along which the same structure is realized at varying intensities, not as a stratum-distinction.} there are four general meanings of \textit{pathos}:
\begin{adjustwidth}{0.5in}{0in}
\begin{enumerate}
\item In its broadest sense, \textit{pathos} ``characterizes a being as something that can in some way be affected by something. Something can \textit{happen} to such a being.'' It is a determination of beings with the character or possibility of ``becoming-otherwise'' (BCAP, p. 131).
\item \textit{Pathos} as the \gk{ἐνέργεια} (\textit{energeiai} ``actualities'') of these alterations: ``the `being-there' of such a shifting occurring-to-one'' (BCAP, p. 131).
\item \textit{Pathos}, narrowing further, ``as the occurring-to-one that has the character of the unpleasant, of the \gk{βλαβερόν} [\textit{blaberon} ``harmful'']. That which happens to me is harmful to me in its happening. This is, indeed, the way that we use the expression `to happen.' But \gk{πάθος} is defined still more precisely: harmfulness is related mostly to \gk{λύπη} [\textit{lypē} ``pain, distress''], so that, as a result, \textit{my attunement to this occurring affects me}. It is a becoming-relevant of something, which aims at my attunement, a becoming-otherwise in the sense of becoming-depressed'' (BCAP, p. 131). The third meaning, here, refines the second: it is not just any occurrence but the occurrence that tones the living being down.
\item Fourth, \textit{pathos} as the magnitude of the harmful. ``\gk{πάθος} designates the `size,' the `measure,' of that which happens to me, that which occurs to me in a harmful way. We have a corresponding expression for that: `that is a blow to me''' (BCAP, p. 132). The fourth meaning narrows still further: it is not just the harmful occurrence but the magnitude of harm --- the gravity of the calamity that strikes.
\end{enumerate}
\end{adjustwidth}

The progression from one sense to the next is not arbitrary; it moves from the most general ontological structure (alterability) through the actual occurrence of alteration (\textit{energeia} of alteration), to the qualitative determination of that occurrence as harmful, to the magnitude of the harm. Heidegger's phenomenological assessment gathers the four under a single encompassing thesis: ``From these four meanings, the \textit{genuine relatedness} of \gk{πάθος} becomes visible; it is related to the \textit{being of living things}, which is characterized by a thus-finding-oneself-again-and-again. The occurring to one befalls and strikes one in this disposition'' (BCAP, p. 132). Taken together, the four senses articulate \textit{pathos} as the structure of \textit{being-occurrable-to} that belongs to the way of being of a finite living being. To live, in the Aristotelian sense Heidegger recovers, is to be ontologically structured as that to which something can happen --- and to which, in concrete factical existence, this and that, again and again, does happen. The occurrence, accordingly, need not without qualification have the character of the harmful; rather, ``Aristotle recognizes a \gk{μεταβολή} [\textit{metabolē} ``change, alteration''], \gk{κίνησις} [\textit{kinēsis} ``movement''], \gk{ἀλλοίωσις} [\textit{alloiōsis} ``alteration, qualitative change''], in which \gk{πάσχειν} [\textit{paschein} ``being-affected''] has the character of \gk{σωτηρία} [\textit{sōtēria} ``preservation'']'' (BCAP, p. 132).\footnote{Aristotle himself: ``Also the expression `to be acted upon' has more than one meaning; it may mean either the extinction of one of two contraries by the other, or the maintenance of what is potential by the agency of what is actual and already like what is acted upon, as actual to potential'' (\textit{De Anima} II.5, 417b2--5).} There is not just one kind of \textit{paschein}: one is destruction by an opposite, the other is preservation by the actuality of what one was potentially. This is the philosophical point Heidegger emphasizes: not all being-affected is loss. Hearing a sound is a being-affected by which my hearing is brought to its full being, not destroyed. Learning is a being-affected by which I come into possession of what I am potentially. \textit{Pathos}, in its widest ontological reach, includes the structure by which a living being is brought into its proper actuality through its encounter with the world.

Each of the four senses of \textit{pathos} is structurally satisfied at the basic affective level of \textit{resonant orexis}. (1) Alterability: \textit{aisthēsis} is constitutionally alterable, since the perceptual \textit{dynamis} is precisely the capacity to be altered by its proper object; perception begins as a \textit{dynamis} that requires an external agent to actualize it. (2) Actual alterations: the \textit{resonant aisthēma} is itself an \textit{energeia} of the \textit{aisthētikon} (perceptive faculty), and \textit{resonant orexis} is the actuality of the \textit{orektikon} (appetitive faculty) at its perceptual-level actualization; Heidegger's gloss --- ``the `being-there' of such a shifting occurring-to-one'' (BCAP, p. 131) --- is just what the affective valence is, the being-there of the soul's being-affected by the pleasant or painful object. (3) Harmful alterations: this moment is seemingly problematic given the narrowing toward \textit{lypē}, but Aristotle's own treatment of perception explicitly includes pain in the normal sensory range\footnote{``Where there is sensation, there is also pleasure and pain'' (\textit{De Anima} II.2, 413b24).} and recognizes that excess is painful or destructive (\textit{De Anima} III.2, 426a30--b8): excess sound deafens, excess light blinds. The harmful-alteration moment is therefore operative at the level of \textit{aisthēsis}. (4) Magnitudes: Aristotle's doctrine that perception has a \textit{logos} or ratio, that excess destroys the ratio, and that stronger stimuli mask weaker stimuli\footnote{``For we must suppose that the stimuli, when equal, tend alike to efface one another, since no one stimulus results from them; while, if they are unequal, the stronger alone is distinctly perceptible'' (\textit{Sense} 7, 447b4--5).} establishes a magnitude-structure at the basic perceptual level. All four senses of \textit{pathos} are accordingly present at \textit{resonant orexis}, and Heidegger's reading of the fourfold as a ``progressive narrowing'' specifies the magnitude-axis along which the same structure is realized at varying intensities and articulations.\footnote{Heidegger's reading of magnitude is quantitative-existential, magnitude as ``size of misfortune.'' The magnitude-axis as deployed in the present synthesis, however, has at least three independent dimensions: hedonic intensity (how pleasant or painful), existential weight (how mattering), and cognitive articulation (how propositionally specified). The articulational dimension carries most of the structural argument and may operate even where hedonic intensity is high --- the drinking case, where real thirst is intense --- without propositional articulation or existential weight in the rhetorical sense.}

\subsubsection{The \textit{Rhetoric}'s Constitutive Threefold: \textit{Pathos} as Change, \textit{Krisis}, and Hedonic Tonality}

The primary definition of emotion in Aristotle's corpus appears not in \textit{De Anima} or any of the traditionally categorized psychological treatises but in the \textit{Rhetoric}: ``the emotions are all those feelings that so change men as to affect their judgements, and that are also attended by pain or pleasure --- such are anger, pity, fear and the like, with their opposites'' (\textit{Rhetoric} II.1, 1378a20--22). Though the definition appears within the context of rhetorical theory's practical applications, it carries substantial ontological weight; and it is this definition along with the four senses of \textit{pathos} that must anchor any systematic analysis of \textit{pathos} in either of its articulational concretions. In his reading of this same passage from the \textit{Rhetoric}, Heidegger isolates the three constitutive aspects of every \textit{pathos}:
\begin{adjustwidth}{0.5in}{0in}
(1) \gk{Μεταβάλλοντες} [\textit{Metaballontes} ``undergoing change'']: something along the way with respect to which ``a change sets in for us,'' through which ``we change'' from one disposition to another. (2) Combine with this change, \gk{διαφέρουσι πρὸς τὰς κρίσεις},\footnote{\textit{diapherousi pros tas kriseis} ``they make a difference to judgments.''} we ``differentiate ourselves'' from ourselves before the change in that which is the hearer's task: ``to take a position,'' ``to form a view.'' The formation of a view involves the manner and mode in which we change. (3) \gk{οἷς ἕπεται λύπη καὶ ἡδονή}:\footnote{\textit{hois hepetai lypē kai hēdonē} ``to which pain or distress and pleasure attend.''} not ``following,'' but rather ``co-given'' in combination with the \gk{πάθη} is a ``being-disposed-as-higher-or-lower'' of the being-there in question. (BCAP, p. 115)
\end{adjustwidth}

\textit{Kriseis}, here, are not external judgments altered by emotion-from-without; they are position-takings of being-there from out of the changed disposition. Heidegger's gloss emphasizes a differentiating-of-oneself-from-oneself, a being-different-from-what-one-was in respect of how one stands toward the matter. Heidegger insists, similarly, that pleasure and pain are not consequences that arrive after the emotion but are co-given with it as a being-disposed-as-higher-or-lower. The \textit{pathos} is the changed-disposition-with-its-position-taking; the pleasure or pain is the tonal elevation or depression co-given with that change. Accordingly, every \textit{pathos}, as articulated in the \textit{Rhetoric} and Heidegger's gloss, is a (1) change-in-disposition that (2) differentiates the subject's position-taking and is (3) co-given with hedonic tonality.

The three-fold articulation is not, however, first introduced at the level of higher-order \textit{pathē}. The second moment in particular --- ``we differentiate ourselves with respect to position-takings'' --- is already operative at the perceptual level on Heidegger's own analysis. He reads Aristotle's explanation of \textit{aisthēsis} as establishing that perception is itself ``a \gk{μέσον} [\textit{mesotēs} ``mean''] with the character of \gk{κριτικόν} [\textit{kritikon} ``discerning''], of the `ability-to-separate' one thing from another\dots It is a being-positioned toward possible objects, which is a \gk{δύναμις} [\textit{dynamis} ``potentiality''] in the sense of \gk{κριτική} [\textit{kritikē} ``faculty of judgment'']'' (BCAP, p. 126). Aristotle confirms this when he argues the perceiving soul is said to make an ``affirmation or negation'' and to pursue or avoid accordingly (\textit{De Anima} III.7, 431a8--12). This basic \textit{resonant orexis} --- the organism's pre-propositional orientation toward the pleasant and away from the painful, what Heidegger names the disclosive \textit{Stimmung} of \textit{Befindlichkeit}\footnote{``What we indicate ontologically by the term `state-of-mind' is ontically the most familiar and everyday sort of thing; our mood, our Being-attuned'' (SZ \S29, H.134; Eng. 172).} --- is co-given with perception from the earliest moment of the actualization chain. It is dispositional; it colors the perceptual field with the most general affective tonality of good or bad, pursue or avoid. The position-taking moment is therefore constitutive of perception with affective valence; it is not introduced anew by \textit{doxa}. \textit{Resonant orexis} satisfies the differentiation aspect pre-propositionally but not pre-positionally --- as the perceiver's already-being-disposed toward the perceived as pursue-able or avoid-able --- whereas \textit{doxa} later re-articulates this position-taking (existential-\textit{hermeneutical} ``as'') propositionally (\textit{apophantical} ``as''). The remaining two aspects similarly hold at the basic level. The first --- \textit{metaballontes} (``those-undergoing change'') --- is satisfied trivially: ``being pained or pleased, or thinking'' are themselves \textit{kineseis} of the soul (\textit{De Anima} I.4, 408b5--7), and basic affective valence is just this kinetic mode at perceptual actualization. The third --- co-givenness of pleasure and pain --- is just what \textit{resonant orexis} is. Emotion (\textit{pathē}), by contrast, is more articulationally concrete. Anger is not undifferentiated displeasure but ``desire accompanied by pain for conspicuous revenge for a conspicuous slight'' (\textit{Rhetoric} II.2, 1378a30). Each emotion, in Aristotle's taxonomy, specifies a propositional evaluation,\footnote{For anger, the slight; for fear, the anticipated evil; for pity, the undeserved suffering (\textit{Rhetoric} II.2, 1378a30--32 for anger; II.5, 1382a21--22 for fear; II.8, 1385b13--15 for pity).} a conative orientation,\footnote{Anger names the conative orientation explicitly --- ``desire accompanied by pain for conspicuous revenge'' (\textit{Rhetoric} II.2, 1378a30). For fear and pity the conative orientation is interpretive, read off the formal structure of avoidance and the wish-that-the-suffering-cease respectively, rather than directly stated by Aristotle in the fear definition (\textit{Rhetoric} II.5, 1382a21--22) or the pity definition (\textit{Rhetoric} II.8, 1385b13--19).} and a hedonic tonality (the pain or pleasure that attends the whole complex). It is this tripartite determinacy that separates emotion from the basic affective valence of pursuit and avoidance. The \textit{Rhetoric}'s three-aspect articulation is therefore not a distinguishing mark of the higher-order \textit{pathē} but the general structure of any \textit{pathos}, applicable at every articulational concretion.
